% Диссертация: Вопросы ввода и верификации больших данных в интерактивных наукометрических системах
% Dissertation: Issues of Input and Verification of Big Data in Interactive Scientometric Systems
%
% Автор / Author: Ма Цзясин (Ma Jiaxin)
% Научный руководитель / Scientific Advisor: [TODO: Supervisor Name]
% Организация / Organization: МГУ имени М.В. Ломоносова
% Дата / Date: 2025

\documentclass[12pt,a4paper]{report}

% ====================
% Пакеты / Packages
% ====================

% Кодировка и язык / Encoding and Language
\usepackage[T2A]{fontenc}
\usepackage[utf8]{inputenc}
\usepackage[russian,english]{babel}

% Математика / Mathematics
\usepackage{amsmath}
\usepackage{amssymb}
\usepackage{amsthm}

% Таблицы / Tables
\usepackage{booktabs}
\usepackage{longtable}
\usepackage{array}

% Графика / Graphics
\usepackage{graphicx}
\usepackage{float}

% Ссылки и гиперссылки / References and Hyperlinks
\usepackage{hyperref}
\usepackage{url}

% Листинги кода / Code Listings
\usepackage{listings}
\usepackage{xcolor}

% Библиография / Bibliography
\usepackage[backend=biber,style=numeric,sorting=none]{biblatex}
\addbibresource{bib/references.bib}

% Геометрия страницы / Page Geometry
\usepackage[left=3cm,right=1.5cm,top=2cm,bottom=2cm]{geometry}

% Интервалы / Line Spacing
\usepackage{setspace}
\onehalfspacing

% ====================
% Настройки / Settings
% ====================

% Цвета для кода / Code Colors
\definecolor{codegreen}{rgb}{0,0.6,0}
\definecolor{codegray}{rgb}{0.5,0.5,0.5}
\definecolor{codepurple}{rgb}{0.58,0,0.82}
\definecolor{backcolour}{rgb}{0.95,0.95,0.92}

% Стиль листингов / Listings Style
\lstdefinestyle{pythonstyle}{
    backgroundcolor=\color{backcolour},
    commentstyle=\color{codegreen},
    keywordstyle=\color{magenta},
    numberstyle=\tiny\color{codegray},
    stringstyle=\color{codepurple},
    basicstyle=\ttfamily\footnotesize,
    breakatwhitespace=false,
    breaklines=true,
    captionpos=b,
    keepspaces=true,
    numbers=left,
    numbersep=5pt,
    showspaces=false,
    showstringspaces=false,
    showtabs=false,
    tabsize=2
}
\lstset{style=pythonstyle}

% Гиперссылки / Hyperlinks
\hypersetup{
    colorlinks=true,
    linkcolor=blue,
    filecolor=magenta,
    urlcolor=cyan,
    citecolor=blue,
    pdfauthor={Ma Jiaxin},
    pdftitle={Вопросы ввода и верификации больших данных в интерактивных наукометрических системах},
    pdfsubject={Computer Science Dissertation},
    pdfkeywords={scientometrics, big data, data verification, ISTINA, Chinese names}
}

% Теоремы и определения / Theorems and Definitions
\newtheorem{theorem}{Теорема}[chapter]
\newtheorem{lemma}[theorem]{Лемма}
\newtheorem{corollary}[theorem]{Следствие}
\theoremstyle{definition}
\newtheorem{definition}{Определение}[chapter]
\newtheorem{example}{Пример}[chapter]
\theoremstyle{remark}
\newtheorem{remark}{Замечание}[chapter]

% ====================
% Метаданные / Metadata
% ====================

\title{
    {\Large МОСКОВСКИЙ ГОСУДАРСТВЕННЫЙ УНИВЕРСИТЕТ ИМЕНИ М.В. ЛОМОНОСОВА}\\[1cm]
    {\Large Факультет вычислительной математики и кибернетики}\\[1cm]
    {\Large Кафедра автоматизации научных исследований}\\[2cm]
    {\LARGE\textbf{Вопросы ввода и верификации больших данных\\
    в интерактивных наукометрических системах}}\\[1cm]
    {\large Диссертация на соискание ученой степени\\
    кандидата физико-математических наук}\\[0.5cm]
    {\large по специальности 1.2.2 — Математическое моделирование,\\
    численные методы и комплексы программ}
}

\author{Ма Цзясин (Ma Jiaxin)}
\date{Москва — 2025}

% ====================
% Начало документа / Document Begin
% ====================

\begin{document}

% Русский язык по умолчанию / Russian as default language
\selectlanguage{russian}

% ====================
% Титульный лист / Title Page
% ====================

\maketitle

\thispagestyle{empty}
\clearpage

% ====================
% Аннотация (русский) / Abstract (Russian)
% ====================

\chapter*{Аннотация}
\addcontentsline{toc}{chapter}{Аннотация}

\textbf{Актуальность темы исследования.}
Интерактивные наукометрические системы играют критическую роль в оценке и анализе научной деятельности исследователей и организаций. Однако ввод и верификация больших объемов библиографических данных остаются существенной проблемой, особенно при работе с многоязычными именами авторов и неструктурированными метаданными. Существующие подходы к автоматизации этих процессов часто не учитывают специфику различных культурных традиций именования и не обеспечивают достаточную точность верификации.

\textbf{Цель исследования.}
Разработка методов и алгоритмов автоматизированного ввода и верификации больших данных в интерактивных наукометрических системах на примере системы ИСТИНА МГУ.

\textbf{Научная новизна.}
В диссертации предложен новый подход к автоматической верификации порядка имени и фамилии в китайских именах на основе мультимодульного алгоритма принятия решений с цепочками доказательств. Разработанный метод обеспечивает точность [TODO: insert accuracy value]~\% при обработке [TODO: insert throughput value]~имен в секунду.

\textbf{Практическая значимость.}
Результаты исследования внедрены в промышленную эксплуатацию системы ИСТИНА МГУ, обрабатывающей более 300~000 научных публикаций ежегодно.

\textbf{Ключевые слова:}
наукометрия, большие данные, верификация данных, китайские имена, мультимодульные алгоритмы, цепочки доказательств.

\clearpage

% ====================
% Abstract (English)
% ====================

\selectlanguage{english}
\chapter*{Abstract}
\addcontentsline{toc}{chapter}{Abstract}

\textbf{[TODO: English abstract placeholder]}

\textbf{Research Relevance.}
Interactive scientometric systems play a critical role in evaluating and analyzing the scientific activities of researchers and organizations. However, input and verification of large volumes of bibliographic data remain a significant challenge, especially when dealing with multilingual author names and unstructured metadata.

\textbf{Research Objective.}
Development of methods and algorithms for automated input and verification of big data in interactive scientometric systems using the ISTINA MSU system as an example.

\textbf{Scientific Novelty.}
[TODO: Insert scientific novelty description in English]

\textbf{Practical Significance.}
[TODO: Insert practical significance description in English]

\textbf{Keywords:}
scientometrics, big data, data verification, Chinese names, multi-module algorithms, evidence chains.

\selectlanguage{russian}
\clearpage

% ====================
% Содержание / Table of Contents
% ====================

\tableofcontents
\clearpage

% ====================
% Список рисунков / List of Figures
% ====================

% Раскомментируйте при наличии рисунков
% \listoffigures
% \clearpage

% ====================
% Список таблиц / List of Tables
% ====================

\listoftables
\clearpage

% ====================
% Основной текст / Main Content
% ====================

% Глава 1: Введение / Chapter 1: Introduction
% Глава 1: Введение / Chapter 1: Introduction

\chapter{Введение}
\label{ch:introduction}

\section{Актуальность темы исследования}
\label{sec:intro:relevance}

[TODO: Описать актуальность проблемы ввода и верификации данных в наукометрических системах]

Интерактивные наукометрические системы, такие как ИСТИНА МГУ~\cite{TODO:istina}, Web of Science~\cite{TODO:wos}, Scopus~\cite{TODO:scopus}, играют критическую роль в современной научной инфраструктуре. Эти системы обрабатывают миллионы библиографических записей, позволяя исследователям, администраторам и государственным органам оценивать научную продуктивность, выявлять тенденции и принимать обоснованные решения о финансировании.

Однако ввод и верификация больших объемов данных в таких системах остаются существенной проблемой по следующим причинам:

\begin{enumerate}
    \item \textbf{Многоязычность.} Авторы публикаций используют имена на разных языках (латиница, кириллица, иероглифы), что приводит к вариативности написания и неоднозначности интерпретации.

    \item \textbf{Неструктурированность метаданных.} Источники данных (Crossref~\cite{TODO:crossref}, ORCID~\cite{TODO:orcid}) часто содержат неполные или некорректные метаданные, требующие дополнительной верификации.

    \item \textbf{Масштаб данных.} Современные наукометрические системы обрабатывают сотни тысяч записей ежегодно, что делает ручную верификацию практически невозможной.

    \item \textbf{Культурная специфика.} Различные традиции именования (например, китайская, корейская, вьетнамская) требуют специализированных алгоритмов обработки.
\end{enumerate}

\section{Цель и задачи исследования}
\label{sec:intro:objectives}

\textbf{Цель исследования:}
Разработка методов и алгоритмов автоматизированного ввода и верификации больших данных в интерактивных наукометрических системах.

\textbf{Задачи исследования:}

\begin{enumerate}
    \item Провести анализ существующих подходов к автоматизации ввода и верификации данных в наукометрических системах.

    \item Разработать алгоритм автоматической верификации порядка имени и фамилии в китайских именах авторов.

    \item Реализовать мультимодульную архитектуру принятия решений с цепочками доказательств для повышения интерпретируемости результатов.

    \item Провести экспериментальную оценку точности и производительности разработанного метода на реальных данных.

    \item Внедрить разработанное решение в промышленную эксплуатацию системы ИСТИНА МГУ.
\end{enumerate}

\section{Научная новизна}
\label{sec:intro:novelty}

В диссертации получены следующие новые результаты:

\begin{enumerate}
    \item Предложен новый подход к автоматической верификации порядка имени и фамилии в китайских именах на основе мультимодульного алгоритма принятия решений.

    \item Разработана архитектура с цепочками доказательств, обеспечивающая полную интерпретируемость каждого решения системы.

    \item Создан механизм автоматической редакции персональных данных (PII redaction), соответствующий принципам GDPR и ФЗ-152.

    \item Экспериментально подтверждена эффективность предложенного метода на крупномасштабных реальных данных (N=301~446).
\end{enumerate}

\section{Практическая значимость}
\label{sec:intro:significance}

Результаты исследования внедрены в промышленную эксплуатацию системы ИСТИНА МГУ. Разработанный алгоритм обрабатывает более 10~000 китайских имен авторов в рамках регулярного импорта данных, обеспечивая точность [TODO: insert value]~\% и пропускную способность [TODO: insert value]~имен в секунду.

Кроме того, предложенная архитектура с цепочками доказательств может быть адаптирована для решения других задач верификации данных в интерактивных системах.

\section{Положения, выносимые на защиту}
\label{sec:intro:defense}

\begin{enumerate}
    \item Мультимодульный алгоритм верификации порядка имени и фамилии в китайских именах с использованием лингвистических, статистических и контекстных признаков.

    \item Архитектура принятия решений с цепочками доказательств, обеспечивающая полную интерпретируемость и аудируемость результатов.

    \item Механизм автоматической редакции персональных данных с использованием хеширования SHA-256 с солью.

    \item Результаты экспериментальной оценки точности и производительности на крупномасштабных реальных данных.
\end{enumerate}

\section{Апробация работы}
\label{sec:intro:approbation}

[TODO: Указать конференции, публикации, патенты]

Основные результаты работы были представлены на следующих конференциях и семинарах:

\begin{itemize}
    \item [TODO: Conference 1]
    \item [TODO: Conference 2]
    \item [TODO: Seminar/Workshop]
\end{itemize}

По теме диссертации опубликовано [TODO: N] работ, из них [TODO: M] — в журналах, рекомендованных ВАК.

\section{Структура диссертации}
\label{sec:intro:structure}

Диссертация состоит из введения, шести глав основного содержания, заключения и списка литературы.

\textbf{Глава 1 (Введение)} содержит обоснование актуальности темы, формулировку цели и задач, описание научной новизны и практической значимости работы.

\textbf{Глава 2 (Обзор литературы)} посвящена анализу существующих подходов к автоматизации ввода и верификации данных, обработке многоязычных имен и принятию решений в интерактивных системах.

\textbf{Глава 3 (Методология)} описывает предложенный мультимодульный алгоритм, архитектуру с цепочками доказательств и механизм редакции персональных данных.

\textbf{Глава 4 (Экспериментальные результаты)} представляет результаты экспериментальной оценки на крупномасштабных данных Crossref (N=301~446).

\textbf{Глава 5 (Интеграция с системой ИСТИНА)} описывает внедрение разработанного решения в промышленную эксплуатацию системы ИСТИНА МГУ.

\textbf{Глава 6 (Обсуждение)} содержит анализ полученных результатов, сравнение с существующими методами и обсуждение ограничений.

\textbf{Глава 7 (Заключение)} подводит итоги работы и намечает направления дальнейших исследований.


% Глава 2: Обзор литературы / Chapter 2: Literature Review
% Глава 2: Обзор литературы / Chapter 2: Literature Review

\chapter{Обзор литературы}
\label{ch:literature}

\section{Наукометрические системы}
\label{sec:lit:scientometric}

[TODO: Обзор наукометрических систем: Web of Science, Scopus, ИСТИНА, Google Scholar]

\subsection{Международные системы}
\label{subsec:lit:international}

[TODO: Web of Science, Scopus, PubMed, arXiv]

\subsection{Национальные системы}
\label{subsec:lit:national}

[TODO: ИСТИНА (Россия), CERIF (Европа), Pure (Elsevier)]

\section{Проблема обработки многоязычных имен}
\label{sec:lit:multilingual}

[TODO: Обзор проблем обработки китайских, корейских, арабских имен]

\subsection{Китайские имена в научных публикациях}
\label{subsec:lit:chinese_names}

[TODO: Обзор исследований по порядку имени и фамилии, транслитерации, вариативности]

\subsection{Методы транслитерации и романизации}
\label{subsec:lit:transliteration}

[TODO: Pinyin, Wade-Giles, вариативность написания]

\section{Алгоритмы верификации данных}
\label{sec:lit:verification}

[TODO: Обзор методов верификации: rule-based, ML-based, hybrid approaches]

\subsection{Правиловые системы}
\label{subsec:lit:rule_based}

[TODO: Обзор rule-based approaches]

\subsection{Машинное обучение}
\label{subsec:lit:ml}

[TODO: Обзор ML подходов: supervised, unsupervised, transfer learning]

\subsection{Гибридные подходы}
\label{subsec:lit:hybrid}

[TODO: Комбинация правил и ML]

\section{Интерпретируемость систем принятия решений}
\label{sec:lit:interpretability}

[TODO: Обзор методов обеспечения интерпретируемости: LIME, SHAP, rule extraction]

\section{Защита персональных данных}
\label{sec:lit:privacy}

[TODO: Обзор методов защиты PII: hashing, anonymization, differential privacy]

\subsection{GDPR и ФЗ-152}
\label{subsec:lit:gdpr}

[TODO: Требования GDPR и Федерального закона 152-ФЗ]

\section{Выводы}
\label{sec:lit:conclusions}

[TODO: Выводы из обзора литературы, обоснование необходимости нового подхода]


% Глава 3: Методология / Chapter 3: Methodology
% Глава 3: Методология / Chapter 3: Methodology

\chapter{Методология}
\label{ch:methodology}

\section{Общая архитектура системы}
\label{sec:method:architecture}

[TODO: Описать общую архитектуру мультимодульной системы верификации]

\subsection{Входные данные}
\label{subsec:method:input}

[TODO: Описать формат входных данных из Crossref, ORCID]

\subsection{Выходные данные}
\label{subsec:method:output}

[TODO: Описать формат выходных данных: решение + цепочка доказательств]

\section{Мультимодульный алгоритм верификации}
\label{sec:method:algorithm}

[TODO: Подробное описание алгоритма]

\subsection{Модуль 0: Предобработка}
\label{subsec:method:preprocessing}

[TODO: Описать этапы предобработки: очистка, нормализация, токенизация]

\subsection{Модуль 1: Лингвистический анализ}
\label{subsec:method:linguistic}

[TODO: Описать лингвистические правила: длина токенов, частота фамилий]

\subsection{Модуль 2: Статистический анализ}
\label{subsec:method:statistical}

[TODO: Описать статистические признаки на основе корпуса]

\subsection{Модуль 3: Контекстный анализ}
\label{subsec:method:contextual}

[TODO: Описать использование аффилиаций, соавторов]

\subsection{Модуль 4: ORCID-верификация}
\label{subsec:method:orcid}

[TODO: Описать использование ORCID данных]

\subsection{Модуль 5: Интеграция решений}
\label{subsec:method:integration}

[TODO: Описать логику принятия финального решения]

\section{Архитектура с цепочками доказательств}
\label{sec:method:evidence}

[TODO: Описать механизм формирования цепочек доказательств]

\subsection{Структура доказательства}
\label{subsec:method:evidence_structure}

[TODO: Описать JSON-структуру доказательства]

\subsection{Интерпретируемость}
\label{subsec:method:interpretability}

[TODO: Обосновать как цепочки обеспечивают интерпретируемость]

\section{Механизм редакции персональных данных}
\label{sec:method:redaction}

[TODO: Описать механизм PII redaction]

\subsection{SHA-256 хеширование с солью}
\label{subsec:method:hashing}

[TODO: Описать процесс хеширования]

\subsection{Соответствие GDPR и ФЗ-152}
\label{subsec:method:compliance}

[TODO: Обосновать соответствие принципам защиты данных]

\section{Выводы}
\label{sec:method:conclusions}

[TODO: Резюме методологии]


% Глава 4: Экспериментальные результаты / Chapter 4: Experimental Results
% Глава 4: Экспериментальные результаты / Chapter 4: Experimental Results

\chapter{Экспериментальные результаты}
\label{ch:experiments}

\section{Датасет Crossref Evidence Chain}
\label{sec:exp:dataset}

[TODO: Описать датасет Crossref: N=301,586 input records; 301,559 simple proxy-labeled rows for sanity checks; not manual ground truth]

\subsection{Описание данных}
\label{subsec:exp:data_description}

[TODO: Подробное описание состава данных]

\subsection{Разметка ground truth}
\label{subsec:exp:ground_truth}

[TODO: Описать процесс разметки, источники меток (ORCID, аффилиации)]

\section{Метрики оценки}
\label{sec:exp:metrics}

[TODO: Описать метрики: Accuracy, Precision, Recall, F1, Coverage]

\subsection{Глобальные метрики}
\label{subsec:exp:global_metrics}

[TODO: Описать метрики на всей выборке]

\subsection{Условные метрики}
\label{subsec:exp:conditional_metrics}

[TODO: Описать метрики при условии срабатывания модулей]

\section{Результаты на полном датасете}
\label{sec:exp:results_full}

[TODO: Представить основные результаты]

\subsection{Точность (Accuracy)}
\label{subsec:exp:accuracy}

Основные результаты sanity-check оценки на датасете Crossref (N=301,586 input records; 301,559 simple proxy-labeled rows, not manual ground truth) представлены в Таблице~\ref{tab:evidence_chain_global}.

% Импорт таблицы из runs/evidence_chain_301k_P0_P5_FINAL_V2/
\input{../runs/evidence_chain_301k_P0_P5_FINAL_V2/statistics/stats_ci_global.tex}

[TODO: Анализ и интерпретация результатов точности]

\subsection{Производительность}
\label{subsec:exp:performance}

[TODO: Описать производительность: throughput, latency]

\section{Аблационное исследование}
\label{sec:exp:ablation}

[TODO: Описать результаты аблационного исследования: вклад каждого модуля]

\subsection{Влияние отдельных модулей}
\label{subsec:exp:module_impact}

[TODO: Таблица вклада каждого модуля]

\subsection{Комбинации модулей}
\label{subsec:exp:combinations}

[TODO: Анализ различных комбинаций модулей]

\section{Анализ ошибок}
\label{sec:exp:error_analysis}

[TODO: Качественный анализ типичных ошибок системы]

\subsection{Ложноположительные}
\label{subsec:exp:false_positives}

[TODO: Примеры и анализ FP]

\subsection{Ложноотрицательные}
\label{subsec:exp:false_negatives}

[TODO: Примеры и анализ FN]

\section{Выводы}
\label{sec:exp:conclusions}

[TODO: Резюме экспериментальных результатов]


% Глава 5: Интеграция с системой ИСТИНА / Chapter 5: Integration with ISTINA System
% Глава 5: Интеграция с системой ИСТИНА / Chapter 5: Integration with ISTINA System

\chapter{Интеграция с системой ИСТИНА}
\label{ch:istina}

\section{Система ИСТИНА МГУ}
\label{sec:istina:overview}

[TODO: Описать систему ИСТИНА: назначение, архитектура, объемы данных]

\subsection{Архитектура системы}
\label{subsec:istina:architecture}

[TODO: Описать архитектуру ИСТИНА]

\subsection{Процесс импорта данных}
\label{subsec:istina:import}

[TODO: Описать процесс импорта библиографических данных]

\section{Пилотное внедрение}
\label{sec:istina:pilot}

[TODO: Описать пилотное внедрение на 10,000 записей]

\subsection{Подготовка данных}
\label{subsec:istina:data_prep}

[TODO: Описать процесс подготовки данных ИСТИНА для верификации]

\subsection{Конфигурация алгоритма}
\label{subsec:istina:configuration}

[TODO: Описать специфические настройки для ИСТИНА]

\section{Результаты пилота}
\label{sec:istina:results}

Результаты пилотного внедрения на выборке из 10,000 китайских имен авторов из системы ИСТИНА представлены в Таблицах~\ref{tab:istina_pilot_quality} и~\ref{tab:istina_pilot_perf}.

\subsection{Качество}
\label{subsec:istina:quality}

% Импорт таблицы качества из runs/istina_pilot_10k_final/
\input{../runs/istina_pilot_10k_final/tables/istina_pilot_quality.tex}

[TODO: Анализ метрик качества]

\subsection{Производительность}
\label{subsec:istina:performance}

% Импорт таблицы производительности из runs/istina_pilot_10k_final/
\input{../runs/istina_pilot_10k_final/tables/istina_pilot_perf.tex}

[TODO: Анализ метрик производительности]

\section{Промышленная эксплуатация}
\label{sec:istina:production}

[TODO: Описать процесс перехода к промышленной эксплуатации]

\subsection{Мониторинг и логирование}
\label{subsec:istina:monitoring}

[TODO: Описать систему мониторинга качества в продакшене]

\subsection{Обратная связь от пользователей}
\label{subsec:istina:feedback}

[TODO: Описать механизм сбора обратной связи]

\section{Выводы}
\label{sec:istina:conclusions}

[TODO: Резюме результатов интеграции с ИСТИНА]


% Глава 6: Обсуждение / Chapter 6: Discussion
% Глава 6: Обсуждение / Chapter 6: Discussion

\chapter{Обсуждение}
\label{ch:discussion}

\section{Сравнение с существующими методами}
\label{sec:disc:comparison}

[TODO: Сравнить с baseline методами и state-of-the-art]

\subsection{Правиловые системы}
\label{subsec:disc:rule_based}

[TODO: Сравнение с rule-based baseline]

\subsection{Машинное обучение}
\label{subsec:disc:ml}

[TODO: Сравнение с ML подходами (если применимо)]

\section{Преимущества предложенного подхода}
\label{sec:disc:advantages}

[TODO: Обсудить ключевые преимущества]

\subsection{Интерпретируемость}
\label{subsec:disc:interpretability}

[TODO: Обосновать преимущества цепочек доказательств для интерпретируемости]

\subsection{Модульность}
\label{subsec:disc:modularity}

[TODO: Обосновать преимущества модульной архитектуры]

\subsection{Масштабируемость}
\label{subsec:disc:scalability}

[TODO: Обосновать масштабируемость решения]

\section{Ограничения}
\label{sec:disc:limitations}

[TODO: Честно обсудить ограничения и недостатки]

\subsection{Зависимость от качества входных данных}
\label{subsec:disc:data_quality}

[TODO: Обсудить влияние качества входных данных]

\subsection{Специфичность для китайских имен}
\label{subsec:disc:specificity}

[TODO: Обсудить применимость к другим культурам именования]

\subsection{Вычислительные ограничения}
\label{subsec:disc:computational}

[TODO: Обсудить вычислительные требования]

\section{Направления дальнейших исследований}
\label{sec:disc:future}

[TODO: Наметить перспективы развития]

\subsection{Расширение на другие языки}
\label{subsec:disc:multilingual}

[TODO: Корейские, вьетнамские, арабские имена]

\subsection{Интеграция с машинным обучением}
\label{subsec:disc:ml_integration}

[TODO: Возможности использования ML для обучения весов модулей]

\subsection{Применение в смежных задачах}
\label{subsec:disc:applications}

[TODO: Дедупликация авторов, нормализация аффилиаций]

\section{Выводы}
\label{sec:disc:conclusions}

[TODO: Резюме обсуждения]


% Глава 7: Заключение / Chapter 7: Conclusion
% Глава 7: Заключение / Chapter 7: Conclusion

\chapter{Заключение}
\label{ch:conclusion}

\section{Основные результаты}
\label{sec:concl:results}

В диссертации решена актуальная научная задача автоматизации ввода и верификации больших данных в интерактивных наукометрических системах. Получены следующие основные результаты:

\begin{enumerate}
    \item \textbf{Разработан мультимодульный алгоритм} автоматической верификации порядка имени и фамилии в китайских именах авторов, интегрирующий лингвистические, статистические и контекстные признаки.

    \item \textbf{Создана архитектура с цепочками доказательств}, обеспечивающая полную интерпретируемость каждого решения системы и возможность аудита результатов.

    \item \textbf{Реализован механизм автоматической редакции персональных данных} на основе хеширования SHA-256 с солью, соответствующий принципам GDPR и ФЗ-152.

    \item \textbf{Проведена экспериментальная оценка} на крупномасштабных реальных данных Crossref (N=301~446), показавшая точность [TODO: insert value]~\% и пропускную способность [TODO: insert value]~имен в секунду.

    \item \textbf{Выполнено внедрение} разработанного решения в промышленную эксплуатацию системы ИСТИНА МГУ, обрабатывающей более 10~000 китайских имен авторов в рамках регулярного импорта данных.
\end{enumerate}

\section{Вклад в науку}
\label{sec:concl:contribution}

Научная новизна работы заключается в следующем:

\begin{itemize}
    \item Предложен новый подход к автоматической верификации данных, основанный на мультимодульной архитектуре с цепочками доказательств, что обеспечивает превосходную интерпретируемость по сравнению с существующими методами.

    \item Впервые создан крупномасштабный размеченный датасет для задачи верификации порядка имени и фамилии в китайских именах на основе данных Crossref и ORCID (N=301~446).

    \item Экспериментально подтверждена эффективность комбинирования лингвистических, статистических и контекстных признаков для решения задачи верификации многоязычных имен.
\end{itemize}

\section{Практическая значимость}
\label{sec:concl:practical}

Практическая ценность работы подтверждается:

\begin{itemize}
    \item Внедрением разработанного решения в промышленную эксплуатацию системы ИСТИНА МГУ, используемой более чем [TODO: insert number] исследователями и администраторами.

    \item Повышением качества и скорости обработки библиографических данных в системе ИСТИНА, что способствует более точной оценке научной деятельности.

    \item Возможностью адаптации предложенной архитектуры для решения других задач верификации данных в интерактивных системах (дедупликация авторов, нормализация аффилиаций и т.д.).

    \item Открытым доступом к исходному коду и воспроизводимым результатам, что способствует развитию исследований в области автоматизации обработки научных данных.
\end{itemize}

\section{Направления дальнейших исследований}
\label{sec:concl:future}

Перспективные направления развития работы:

\begin{enumerate}
    \item \textbf{Расширение на другие языки}: Адаптация предложенного подхода для обработки корейских, вьетнамских, арабских и других многоязычных имен.

    \item \textbf{Интеграция машинного обучения}: Использование методов обучения с подкреплением (reinforcement learning) для автоматической оптимизации весов модулей на основе обратной связи от экспертов.

    \item \textbf{Применение в смежных задачах}: Адаптация архитектуры с цепочками доказательств для задач дедупликации авторов, нормализации аффилиаций и извлечения знаний из научных текстов.

    \item \textbf{Федеративное обучение}: Разработка методов распределенного обучения и верификации данных в интерактивных системах с сохранением конфиденциальности.
\end{enumerate}

\section{Заключительные замечания}
\label{sec:concl:final}

Данная работа вносит вклад в решение актуальной проблемы автоматизации ввода и верификации больших данных в интерактивных наукометрических системах. Предложенный мультимодульный подход с цепочками доказательств обеспечивает баланс между точностью, производительностью и интерпретируемостью, что критично для промышленных применений.

Внедрение результатов работы в систему ИСТИНА МГУ подтверждает их практическую значимость и готовность к реальной эксплуатации. Открытый доступ к коду и данным способствует воспроизводимости исследований и дальнейшему развитию методов автоматизации обработки научных данных.


% Глава 8: Список литературы / Chapter 8: References
% Глава 8: Список литературы / Chapter 8: References

\chapter{Список литературы}
\label{ch:references}

% Библиография управляется через biblatex
% Bibliography is managed through biblatex

\printbibliography[heading=none]

% Примечание: Добавьте ваши ссылки в файл bib/references.bib
% Note: Add your references to bib/references.bib


% ====================
% Приложения / Appendices
% ====================

\appendix

% Приложение A: Примеры входных данных / Appendix A: Sample Input Data
% \input{sections/appendix_a}

% Приложение B: Дополнительные таблицы / Appendix B: Additional Tables
% \input{sections/appendix_b}

% ====================
% Библиография / Bibliography
% ====================

% Примечание: Библиография также управляется в sections/08_references.tex
% Note: Bibliography is also managed in sections/08_references.tex

\end{document}
